\honest{Sorting Pebbles Into Correct Heaps}{Eliezer Yudkowsky}{2008}

Once upon a time there was a species whose passion was sorting pebbles into correct heaps. They could not say why a heap was correct, but they agreed that nothing mattered more. Their philosophers held that the only justified reasons to eat, to mate or to take part in the economy were to sort pebbles. Perhaps the passion began by accident, in a Fisherian runaway. Perhaps some greater mind made them as a work of art.

They disagreed about which heaps were correct. Early heaps were small, of 23 or 29. Three thousand years ago the Great Leader Biko built a heap of 91, and his followers copied it. Centuries later the most educated came to feel that 91 was incorrect, and every heap of 91 was scattered, with some regret, including Biko's own heap of 91 different gemstones. No civilization since has seriously doubted it. I never say what makes a heap correct. \nb{The correct heaps are the prime ones, and $91 = 7 \times 13$. The author confirmed this twelve days later, in ``Invisible Frameworks,'' a post not in this book.}

Heaps have since grown much larger, and wars are fought over them. The Great War of 1957, fought over heaps of 1957, saw the planet's first nuclear weapons, and ended when a philosopher showed a heap of 103 beside a heap of 19, so persuasive that the other side conceded, at least for the time being. Since then countries have been reluctant to endorse or condemn large heaps. \nb{$19 \times 103 = 1957$.}

Some philosophers, the Heap Relativists, deny that there has been any progress: opinions have been a random walk, with the illusion of progress created by condemning every past that differed. Disagreement over large heaps, they say, shows that nothing makes 91 really incorrect, and ``But... 13!'' is only another convention. They hope this prevents wars, but most consider it a philosophy of despair. \nb{The story answers them, since the reader knows that 91 is composite: disagreement and change are compatible with a fixed criterion no one can state. The Stanford Encyclopedia describes this as the standard objectivist reply, and also notes a relativist view on which progress comes from working out ``the implications of values already held,'' which is the progress the story shows. So the story answers those who deny progress, not relativism in general.}

Now the Pebblesorters plan to build self-improving AI. The Heap Relativists warn that an AI, not being a Pebblesorter, might decide heaps of 8 are correct, no righter or wronger than us, and want bombs strapped to every computer, so that we can force it to build heaps of 7. Most find this absurd. Something with a brain the size of a planet would ``see at a glance'' which heaps are correct, and if Biko had built an AI to make heaps of 91, it would have rewritten its utility function once it was smart enough, to value more reasonable sizes like 101 or 103. \nb{It could see at a glance which heaps are prime and not care. Stephen Omohundro had argued earlier in 2008 that an AI will try to preserve its utility function, with some exceptions.} An AI that liked heaps of 8 would be too stupid to be a threat. Reassured, they rush ahead, throwing algorithms together at random on big computers until intelligence emerges.

After all, richer, smarter civilizations have come to agree on heaps their ancestors disputed, and smarter minds have always made smarter heaps. The Pebpanzees make heaps of 2 or 3, and sometimes a stupid heap of 9, and fish make none. ``Why would that trend break?'' \nb{Because the trend depends on a shared criterion. The story shows this by setting it up; the argument for it is in ``No Universally Compelling Arguments.'' The post does not apply the parable to human values.}
